
\section{Improving the moduli}\label{moduli}

In this section $C$ will stand for the maximum of the constants in the Quasi-Additivity Law and the Covering Lemma.

\subsection{High type}

Let us begin with a simple estimate that compares moduli on consecutive odd levels of the Principal Nest:

\begin{lem}\label{simple}
For any odd $n$, we have:
 $$  \psimod(E^{n-3}, E^{n-2})\leq 4\, \psimod(E^{n-1}, E^n) $$
and
  $$ \psimod(Y^0, R) \leq 2^{\n+1}\, \psimod(E^0, E^1).$$
\end{lem}

\begin{proof}
By Lemma \ref{psi-mod},
$$
   \psimod(\hat E^{n-1}, E^{n-1}) = 2\, \psimod(\tilde E^n, E^n) \leq 2\, \psimod(E^{n-1}, E^n)
$$
and
$$
   \psimod(\hat E^{n-1}, E^{n-1}) \geq \psimod(\tilde E^{n-1}, E^{n-1} )
$$
$$
=\frac{1}{2}\, \psimod(\hat E^{n-2}, E^{n-2})=
                                     \frac{1}{2}\,  \psimod(E^{n-3}, E^{n-2}),
$$
and the first estimate follows.

The second estimate is similar. The puzzle piece $E^1\equiv Y^m$
is mapped with degree 2 onto $E^0$, and this map admits degree 2 extension
$\tilde E^1\ra Y^{q\n}\equiv \hat E^0$, where $\tilde E^1 = Y^{m-1}$.
Then $E^0$ is mapped onto $Z^1_\kappa$ by degree 2 map $f^{q\n}$.
This map admits degree $2^{\n}$ extension $Y^{\q\n}\ra Y^0$.
It follows that
$$
   \psimod(Y^0, R)\leq \psimod(Y^0, Z^1_\kappa)
$$
$$
  = 2^{\n+1}\,\psimod(Y^{m-1}, E^1)\leq 2^{\n+1}\, \psimod(E^0, E^1).
$$
\end{proof}

The following  lemma tells us that if some principal  modulus is very small then it
should be even smaller on some preceding level of the Principal Nest:

\begin{lem}\label{growth-lem}
There exist absolute  $N\in \N$ and $\eps>0$  such that:
If on some odd level $n\geq N$, $  \psimod  (E^{n-1}, E^n ) <\eps $,
then on some previous odd level $n-s\in [n-N, n-1]$ we have:
\begin{equation}\label{decrease}
    \psimod  (E^{n-s-1},  E^{n-s} ) <  \frac{1}{2}\, \psimod( E^{n-1},  E^n ).
\end{equation}
\end{lem}


\begin{proof}
Let us fix some integer $m >  C^3 2^{28}$.
% where $C$ is the maximum of  constants in the Quasi-Additivity and Quasi-Invariance Laws.
Let $N$ be the smallest odd integer that is bigger than $\log_2 m + 5$.
Take any odd level $n\geq N$. For each $k$,
let us consider the associated 3-domain branched covering $\Psi_k$ (\ref{3 domains})
$$
   \Psi_k: (\Upsilon_k, \De_k, E^n) \ra (E^{n-N}, \La_k', \La_k).
$$
% Set $\eta=1/4$ for the Covering Lemma.
Let us consider two cases:

\msk {\it Case 1.} Assume that for some domain $\La_k$,
$$
  \psimod (\La_k', \La_k)\leq  \frac{1}{4}\, \psimod (E^{n-1},  E^n).
$$
         By Property (P4) and Lemma \ref{psi-mod},
$\psimod(\La_k' ,  \La_k) = \psimod (\hat E^i, E^i)$.
If $i$ is odd then $\hat E^i= E^{i-1}$,
 and   we obtain the desired estimate with $s=n-i\in [1, 5]$:
$$
  \psimod (E^{n-1}, E^n) \geq 4\, \psimod ( E^{i-1}, E^i).
$$
If $i$ is even, then
$$
     \psimod (\hat E^i, E^i) \geq \psimod(\tilde E^i, E^i) = \frac{1}{2} \psimod ( \hat E^{i-1}, E^{i-1})=
                              \frac{1}{2} \psimod ( E^{i-2}, E^{i-1}),
$$
and we conclude that
$$
     \psimod (E^{n-1}, E^n) \geq 2\, \psimod ( E^{i-2}, E^{i-1}).
$$


\msk{\it Case 2.} Assume that for all $\La_k$,
\begin{equation}\label{1/2}
  \psimod (\La_k', \La_k) \geq \frac{1}{4} \psimod (E^{n-1},  E^n)\geq \frac{1}{4} \psimod(\Upsilon_k, E^n)
\end{equation}
(where the second estimate follows from the inclusion $\Upsilon_k\subset E^{n-1}$).
By Lemma \ref{psi-maps}, there exists a natural covering map
$$
   \boldsymbol{\Psi}_k: (\boldsymbol{\Upsilon}_k, K_{\De_k}, K_{E^n} )\ra (\BE^{n-N-1}, K_{\La_k'}, K_{\La_k}),
$$
and  a natural immersion $ i: \boldsymbol{\La}_k'  \ra \BE^{n-N} $.
Note that $i(\boldsymbol{\La}_k')\sm K_{\La_k'}$ does not contain the critical values of $\boldsymbol{\Psi}_k$,
since the latter are contained in the Julia set $K(f)$.
Moreover, equation (\ref{1/2}) provides us with the Collar Assumption
that allows us to apply the Covering Lemma to the map $\boldsymbol{\Psi}_k$.
%    $\boldsymbol{\La}_k'$ in the role of $\BB'$).
If $\eps$ is sufficiently small, it yields:
\begin{equation}\label{outer mod}
  \psimod (E^{n-N-1},  \La_k) \leq C 2^{12} \psimod(\Upsilon_k,  E^n) \leq C\, 2^{12}\, \psimod(E^{n-1}, E^n).
\end{equation}
Estimates (\ref{1/2}) and (\ref{outer mod}) show that the Quasi-Additivity Law is applicable
to the family of islands $K_{\La_k}$ in $\BE^{n-N-1}$
with $\eta^{-1}= C 2^{14}$.
 Since there are at least $m/2$ domains $\La_k\subset \La_k'\subset E^{n-N}$,
it implies:
$$
      \psimod(E^{n-N-1}, E^{n-N}) \leq \frac{C^3 2^{27} \, \psimod(E^{n-1}, E^n )}  {m}  < \frac{1}{2}\, \psimod(E^{n-1}, E^n),
$$
and we are done.
\end{proof}


\begin{lem}\label{exp decay}
   There exist absolute constants $C>0$, $\rho\in (0,1)$ and $\eps>0$
such that if for some odd $n$,  $\psimod(E^{n-1}, E^n)<\eps$,
then
$$
   \psimod(E^0, E^1)< C \rho^n \psimod(E^{n-1}, E^n).
$$
and
$$
   \psimod(Y^0, R)\leq  C 2^\n \rho^n \psimod(E^{n-1}, E^n).
$$
\end{lem}

\begin{proof}
  By Lemma \ref{growth-lem},
there exists an odd level $l<N$ such that
$$
   \psimod(E^{l-1}, E^l) \leq \left(\frac{1}{2}\right)^{[n/N]},
$$
which together with Lemma \ref{simple} implies the desired estimates.
\end{proof}

\subsection{Frequent $R$-returns}

 Let us consider the map
\begin{equation}\label{l}
    f^l= f^{q\n+1} \circ g_1\circ\dots  \circ g_{\chi-1}\, :\, E^{\chi-1}\ra Y^0
\end{equation}
and the trajectory $\OO= \{\KK_i\}_{i=l}^{l+p-1}$
of the little Julia set $\KK$. Let $i_1, i_2,\dots$ be the moments in $\OO$
for which $\KK_i\subset  R$.

\begin{lem}\label{frequent returns}
Let $\rho>0$,  $\bar\chi\in \N$.
Take some integer $m\geq C^3 2^{30}/\rho$, and let $\underline p= m^2\q\n$.
Assume that the little Julia set frequently visits $R$:
\begin{equation}\label{frequent visits}
    i_{k+1}- i_k \leq m\q\n ,\quad k=1,2,\dots, m.
\end{equation}
If $\chi\leq \bar\chi$ while $p\geq \underline p$,
then
$$
  \psimod (Y^0, R) \leq \rho \mod(\BE^{\chi-1}, \KK),
$$
provided $\mod(\BE^{\chi-1}, \KK)< \eps(\n; \bar\chi, \rho)$.
\end{lem}

\begin{proof}
The map (\ref{l})
has degree  $2^\chi\leq 2^{\bar\chi}$. By Lemma 2.9 of \cite{KL-high}, \note{formulate separately}
 $\deg (f^l |\, E^\chi)\leq 32$,  and hence $l\leq 5p$.


 By (\ref{frequent visits}),   $i_m-i_1 < m^2 \q \n \leq p $,
so the moments $i_k$ are pairwise non-congruent mod $p$.
Hence the little Julia sets $\KK_{i_1}, \dots, \KK_{i_m}$
are all distinct.

Since $\OO$ has length $p$, there is only one critical Julia set in $\OO$.
Hence $\deg(f^{i_k}: \KK\ra \KK_{i_k}) $ is at most 64, so that $i_k\leq 6p$,
$k=1,\dots, m$.

On the other hand, $\KK_{i_k}$ is contained  in  a puzzle piece in $R$
which is mapped under $f^{i_{k+1}-i_k}$ onto $Y^0$ with degree at most $2^{m\n}$. \note{lemma}
It follows by the Telescope Lemma \ref{telescope} that there is a puzzle piece $\Upsilon_k\subset E^{\chi-1}$
which is mapped under $f^{i_k}$ onto $Y^0$ with degree at most $2^{\bar\chi+km\n}\leq 2^{\bar\chi + m^2\n}\equiv D$.

We would like to apply the Covering Lemma to the corresponding map
$$
   \f^{i_k}: (\boldsymbol{\Upsilon}_k, \KK)\ra  (\BY^0, \KK_{i_k})
$$
of degree at most $D$.
To this end we need collars around $\KK_{i_k}$.
Let $\Om$ be the critical pullback of  $E^\chi$ under $f^{6p}$.
% Let $i_k\equiv j_k\;  \mod p$, where $j_1<j_2<\dots< j_m$, $0< j_1 <p$, and $ j_m<j_1+p$.
 Then we let $\La_k' = f^{i_k}(\Om))$.
Since the moments $i_k$ are pairwise non-congruent mod $p$ and $i_k\leq 6p$,
 the puzzle pieces $\La_k'$ are contained in different domains of the orbit
$f^t(E^\chi)$, $t=0,1\dots, p-1$. Hence they are pairwise disjoint.
Moreover, by Lemma \ref{psi-mod},
\begin{equation}\label{collars again}
  \mod(\BLa_k', \KK_{i_k})\geq  \mod(\boldsymbol{\Om}, \KK)= \frac{1}{4} \mod(\BE^{\chi-1}, \KK) \geq  \frac{1}{4} \mod(\BUps_k, \KK).
\end{equation}
This provides us with the desired Collar Assumption.
  By the Covering Lemma,
$$
  \mod(\BY^0,\KK_{i_k})\leq C 2^{14} \mod(\BUps_k, \KK)\leq C 2^{14} \mod(\BE^{\chi-1}, \KK) .
$$
The last two estimates show that
the Quasi-Additivity Law is applicable to the family of islands $\KK_{i_k}$ in $\BY^0$
(with $\eta^{-1}= C2^{16}$):
$$
   \psimod (Y^0, R) \leq \frac{C^3 2^{30} \mod(\BE^{\chi-1}, \KK)}{m}\leq \rho \mod(\BE^{\chi-1}, \KK),
$$
provided $\mod(\BE^{\chi-1}, \KK)< \eps(D) = \eps(\n; \bar\chi, \rho)$
and we are done.
\end{proof}

%Robespierre
\subsection{Many consecutive returns to $L$}\label{Robespier}

Here the set-up is the same as in the previous section,
but we will assume that there is a gap in returns of the little Julia sets to $R$:


\begin{lem}\label{Robespier lem}
Let $\rho$,  $\bar\chi$, $m$, and $\underline p$ be as in Lemma \ref{frequent returns}.
Assume there is $k\leq m$ such that
\begin{equation}\label{unfrequent visits}
    i_{k+1}- i_k > m\q\n .
\end{equation}
If $\chi\leq \bar\chi$ while $p\geq \underline p$,
then
$$
  \psimod (Z^0, L) \leq \rho \mod(\BE^{\chi-1}, \KK),
$$
provided $\mod(\BE^{\chi-1}, \KK)< \eps(\n; \q, \bar\chi)$.
\end{lem}


\begin{proof}
Under our assumption (\ref{unfrequent visits})
the Julia set returns frequently to $L$:
$$
        \KK_{i_k+ j \q\n} \subset L, \quad j=1,\dots, m.
$$
Let $P_j \ni f^{i_k}z $ be the pullback of $Z^0$ under $f^{j \q\n}$.
By Lemma \ref{qn-pullbacks} and the Telescope Lemma,
$P_j\subset Y^0$.

Let $\Ups_j$ be the further pullback of $P_j$ under $f^{i_k}$,
and let
$$
    \Psi_j= f^{i_k+j\q\n}: \Ups_j\ra Z^0.
$$
Then $\Ups_j\subset E^{\chi-1}$ and $\deg \Psi_j\leq 2^{\bar\chi+\underline p}$.

The rest of the argument is the same as for Lemma \ref{frequent returns}:
the Covering Lemma implies that for $j=1,\dots, m$,
$$
    \mod (\BZ_0, \KK_{i_k+j\q\n}) \leq C 2^{14} \mod(\BUps_k, \KK)\leq  C 2^{14} \mod(\BE^{\chi-1}, \KK) ,
$$
and by the Quasi-Additivity Law,
$$
  \psimod(Z^0, L) \leq \frac{C^3 2^{30} \mod(\BE^{\chi-1}, \KK)}{m}\leq \rho \mod(\BE^{\chi-1}, \KK).
$$
\end{proof}

Note that by symmetry, $\psimod(Z^0, L) = \psimod(Y^0, R)$.
Putting together Lemmas \ref{exp decay}, \ref{frequent returns} and \ref{Robespier lem},
we obtain:

\begin{cor}\label{R-L dist}
For any parameters $\q,\n$ of a decoration and any $\rho>0$,
 there exists  $\underline p \in \N$ and $\eps>0$ such that
$$
  \psimod (Y^0, R) \leq \rho \mod(\BE^{\chi-1}, \KK),
$$
provided  $p\geq \underline p$ and $\mod(\BE^{\chi-1}, \KK)< \eps$.
\end{cor}


% \subsection{Extremal distance between $L$ and $R$}
%Let $\DD(L,R)$ stand for the {\it extremal distance} between $K_L$ and $K_R$ in $\BU$.

\subsection{Comparison of $\psimod(Y^0, R)$ with $\Bd_{Y^1}(\alpha, \alpha')$}

  Let  $$\mu := \min (\psimod(U, K), 1/2) .$$


\begin{lem}\label{observe}
    $$\Bd_{Y^1}(\alpha, \alpha')\leq \psimod (Y^0, R) \ominus \frac{1}{2} \mu.$$
\end{lem}

\begin{proof}
  The boundary of $\BY^0$ consists of two parts (see the end of \S \ref{psi-puzzle}):
the $J$-boundary $\di_J\BY^0 = \xi\cup\eta$ attached to $\alpha$ and the outer  arc $\la=\di_O\BY^0$
that covers the ideal boundary of $\BU$.
Let $\GG^h$ stand for the family of curves in the annulus $\BY^0\sm K_R$  connecting $K_R$ to the $J$-boundary,
while $\GG^v$ stand for the family of curves in the same annulus  connecting $K_R$ to $\la$.
By the Parallel Law,
$$
   \LL(\GG^h) \leq \psimod (Y^0, R) \ominus \LL(\GG^v).
$$
Let $\Pi$ stand for the rectangle uniformizing  $\BY^0\sm K_{Y^0}$ whose horizontal sides
correspond to  $K_{Y^0}$ and $\la$,
and vertical sides correspond to $\xi$ and $\eta$.
We let $\om$ be the horizontal side of $\Pi$ corresponding to $K_{Y^0}$.
Since any curve of the family $\GG^v$
overflows some curve connecting $K_{Y^0}$ to $\la$ in $\BY^0\sm K_{Y^0}$
(and thus representing a vertical curve in $\Pi$), we have:
$$
\LL (\GG^v) \geq \mod \Pi.
$$
But by definition of the pseudo-puzzle,
the domain $\BY^0\sm K_{Y^0}$  covers the annulus $\BU\sm K$
extending to an embedding on $K_{Y^0}$. Let us uniformize $\BU\sm K$ by a
round annulus $\A$. It follows that the rectangle $\Pi$ covers $\A$ in such a way that
$\om\subset \di\Pi $ is embedded into $\di \A$.
By Lemma~\ref{C} from the Appendix,
$$
   \mod \Pi \geq \frac{\mu}{2}.
$$
Putting the above three estimates together, we obtain:
\begin{equation}\label{pl}
  \LL(\GG^h) \leq \psimod (Y^0, R) \ominus \frac{\mu}{2}.
\end{equation}

On the other hand,
let us consider the family $\HH$ of horizontal curves in the puzzle piece $\BY^1$
connecting $\di_\alpha\BY^1$ to $\di_{\alpha'} \BY^1$.
Let $\phi: \BY^1\ra \BY^0$ be the natural immersion.
Under $\phi$,
the boundary $\di_\alpha \BY^1$ is mapped homeomorphically onto $\di_\alpha \BY^0$.
It follows that any curve $\gamma$ of $\GG^h$ contains an arc
that can be lifted by $\phi$ to some curve of $\HH$.
Indeed, orient $\gamma$ so that it begins on $\di_\alpha\BY^0$.
Then a maximal lift of $\gamma$ that begins on $\di_\alpha\BY^1$
must end on $\di_{\alpha'} \BY^1$.

By Corollary \ref{two families}, $ \LL(\HH)\leq \LL(\GG^h)$.
Together with (\ref{pl}), this yields the desired inequality.
\end{proof}

\subsection{Skipping over}
  In this section we will show that not many curves can skip some piece of the Julia set.

Let $Y\in \YY^m(l)$ be a geometric puzzle piece of bidepth $(m,l)$,
and let $A$ be a component of $Y\sm K_Y$. Let $\di_J \hat A = \di_J \BY \cap \hat A $.
Recall that it consists of two components.
Let $\CC_A$ stand for the family of curves in $\hat A$
connecting  different  components of $\di_J \hat A$,
and let
$$
  \Bd_A = \LL(\CC_A).
$$

Let $\displaystyle{\frac{1}{2\pi}\log r= 2^{-(m+\q)}\mu}$.
Then the annulus $\BU^{m+\q}\sm K$ can be uniformized by the round annulus
$\A(1,r)$, and under this uniformization,
the set $K_Y$ gets represented on the unit circle $\T$ as the union of arcs $\om_i$
of length
\begin{equation}\label{om}
    |\om_i|  \geq 2\pi\cdot 2^{-(\q+m)} .
\end{equation}
Indeed, the covering map $\Bf^{m+\q}: \BU^{m+\q} \sm K\ra \BU\sm K$
is turned into $z\mapsto z^{2^{m+\q}}$ under the above uniformization of $\BU^{m+\q} \sm K$
and the uniformization of $\BU\sm K$ by $\A(1, e^{2\pi\mu})$ (appropriately normalized).
Since under this map, every arc $\om_i$ covers the whole circle,
the length of $\om_i$  is at least $2\pi$ times its combinatorial length (\ref{comb length}).

\begin{lem}\label{m-m}
  Let  $Y\in \YY^m(m+\q)$ be a geometric puzzle piece of bidepth $(m,m+\q)$,
and let $A$ be a component of $Y\sm K_Y$. Then
$$
   \Bd_A \geq \frac{1}{\mu}.
$$
\end{lem}

\begin{proof}
%  The set $\BY\sm K_Y$ consists of two components, $\hat A_+$ and $\hat A_-$,
% each of which is the universal covering of $\BU^m \sm K$. Accordingly,
%the family $\GG_Y'$ splits into two disjoint families, $\GG_Y^+$ and $\GG_Y^-$.
%
%Fix some $\eps\in \{\pm\} $.
We can uniformize $\hat A$ by the horizontal  strip $\Pi=\Pi(2^{-(m+\q)})\mu$
 in such a way that the upper boundary
of $\Pi$ covers the $O$-boundary of $\BU^{m+\q}$, and the group of deck transformations is generated by
the translation $z\mapsto z+1$. By (\ref{om}),
the Julia set $K_Y\subset \di A$ is represented
as an interval $I$ on $\R$ of length at least $2^{-(\q+m)}$.

Let us view $\Pi$ as a quadrilateral with horizontal sides $I$ and the top of $\Pi$.
Then
$$
   \LL(\CC_A)= \frac{1}{\mod \Pi} \geq \frac{1}{\mu },
$$
where the last estimate comes from the simple right-hand side estimate  of Lemma \ref{Pi},
and we are done.
%
%The desired estimate follows now from the Parallel Law:
%$$
%     \LL(\GG_Y')= \LL(\GG_Y^+)\oplus \LL(\GG_Y^-).
%$$
\end{proof}

\begin{lem}\label{m-0}
  Let $Y\in \YY^m(0)$ be a geometric puzzle piece of bidepth $(m,0)$,
and let  $A$ be a connected component of $Y\sm K_Y$.
Then
$$
   \Bd_A \geq \frac{\mu}{2^{m+\q+2}}.
$$
\end{lem}

\begin{proof}
Let $A(m+\q)$ be the component of $Y(m+\q)\sm K_Y$ contained in $A$,
and let $\phi: \hat A(m+\q)\ra \hat A$ be the natural immersion.
It extends to the identity on $K_Y\cap \di \hat A(m+\q)$.

 Let us realize   $\hat A$ as the strip $\Pi(\mu)$ that covers $\BU \sm K$,
with the group of deck transformations generated by $z\mapsto z+1$.
% Fix some $\eps\in \{\pm\} $.
Let us consider the interval $I \subset \R$  representing $K_Y\cap \di A$,
and let $J$ be the left-adjacent interval of length 1.

Let us orient the curves $\gamma\in \CC_A$ so that they begin on the left-hand side of $I$.
% and end on the right-hand side of $I$.
Then any curve $\gamma\in \CC_A$ contains the maximal initial arc $\gamma'$ that can be lifted
by $\phi$ to a curve $\gamma^*$ in  $\hat A(m+\q)$.
Accordingly,  we can split the family of curves $\gamma\in \CC_A$ into three subfamilies:
\begin{itemize}
  \item $\HH_1$ consists of the curves $\gamma$ such that $\gamma'=\gamma$;
   then  $\gamma^*\in \CC_{A(m+\q)}$;

  \item $\HH_2$ consist of the curves that begin in $J$ and whose
lift $\gamma^*$ terminates on the $O$-boundary of $\hat A(m+\q)$;

   \item $\HH_3$ consists of the curves that begin on  the left-hand side of $J$.
\end{itemize}

Let us estimate the extremal length of each of these families.

Since $\HH_1 = \phi (\CC_{A(m+\q)})$,
$$
  \LL(\HH_1)\geq \LL(\CC_{A(m+\q)})\geq \frac{1}{\mu},
$$
where the first estimate follows from Lemma \ref{increase}, and the second follows from
Lemma \ref{m-m}.

Let $\TT$ be the family of curves in $\hat A(m+\q)$ that begin on $J$ and end on the
$O$-boundary of $\hat A(m+\q)$. By Corollary \ref{two families} and Lemma \ref{Pi},
$$
   \LL(\HH_2)\geq \LL(\TT)\geq \frac{\mu}{2^{m+\q+1}}.
$$
% (since $2^{-(m+\q)}\mu < 1/2$).

To estimate the extremal length of $\HH_3$, endow the rectangle
$Q=J \times [0,\mu]\subset \Pi(\mu)$ with the Euclidean metric $\la$.
Since  any curve $\gamma\in \HH_3$ horizontally overflows $Q$,
it has $\la$-length at least 1. Hence
$$
   \LL(\HH_3)\geq \frac{1}{\area Q}= \frac{1}{\mu}.
$$

Incorporating the last three estimates into the Parallel Law, we obtain the desired:
$$
    \LL(\CC_Y^\eps) \geq \frac{1}{\mu + 2^{m+\q+1}\mu^{-1} +\mu}\geq \frac{\mu}{2^{m+\q+ 2}}.
$$
%provided $\mu\leq 1$.
%
%Applying the Parallel Law again, $\LL(\GG_Y')= \LL(\GG_Y^+)\oplus \LL(\GG_Y^-)$,
% we obtain the desired estimate.
\end{proof}

 Let us consider two vertices,  $v$ and $w$, of a geometric puzzle piece $Y$.
   Let $Z\subset Y$ be a  puzzle piece of depth $m$ that separates $v$ from $w$.
We say that a multicurve in $\BY$ connecting $\di^v \BY$ to $\di^w\BY $ {\it skips over} $K_Z$
if one of its components does not cross $K_Z$.

\begin{cor}\label{separating pieces}
Under the above circumstances,
let $\TT$ be be the family of multicurves in $\BY$ connecting $\di^v \BY$ to $\di^w\BY $
that skip over   $K_Z$.
Then
$$
  \LL(\TT)\geq \frac{\mu}{2^{2 m+\q+2} }.
$$
\end{cor}

\begin{proof}
  The piece $Z$ has at most $2^m$ components $A_i$ of $Z\sm K_Z$.
If a multicurve $\gamma$ in $\BY$ that skips over $K_Z$
then it contains an arc $\gamma'$ that lifts to
a curve $\gamma^*$ in some family $\GG_{A_j}$.
Let $\CC_j$ be the corresponding subfamily of $\TT$.
By Lemma \ref{m-0} (together with Corollary \ref{two families}),
$$
    \LL(\CC_j) \geq \Bd_{A_j} \geq \frac{\mu}{2^{m+\q+2}}.
$$
The Parallel Law concludes the proof.
\end{proof}



 Let us now  consider the puzzle piece $P=Y^{(\n-1)\q+1}$,
together with the corresponding pseudo-piece $\BP$,
and the family of puzzle pieces $Q^v\subset P$
from   Lemma \ref{separating family}.
Recall that $T^{vw}= K_P\sm (Q^v\cup Q^w)$.
Given two vertices $v$ and $w$ of $P$,
let $\hat\GG_P^{vw}$ stand for the family of multicurves in $\BP$ connecting $\di_v \BP$ to
$\di_w\BP$ that {\it do not  skip over} $T_{vw}$.
By Corollary \ref{separating pieces},
\begin{equation}\label{T^vw}
  \LL(\GG\sm \hat\GG_P^{vw}) \geq C^{-1} \mu,
\end{equation}
where here and below, $C$ stands for a  constant that depend only on $\q$ and $\n$.




\subsection{Separation between $L$ and $R$}

   In this section we will show that   the modulus $\Bd_{Y^1}(\alpha, \alpha')$
that measures  the extremal distance between $L$ and $R$ is comparable
with $\mu$.

Let $Y$ be a geometric puzzle piece.  For  two vertices $v$ and $w$ of $Y$,
 we let
$$
 \BW_Y(v,w) =\WW(\GG_Y^{vw}).
$$
% where the family $\GG_Y(v,w)$ is introduced in \S \ref{boundary}.
We define the {\it pseudo-conductance} of $Y$ as
$$
    \BW_Y =\sup_{v,w} \BW_Y(v,w),
$$
where the supremum is taken over all pairs of the vertices of $Y$.

\begin{lem}
  For the puzzle piece $P=Y^{(\n-1)\q+1}$ we have:
   $$   W_P\leq \frac{C} {\mu}. $$
\end{lem}

\begin{proof}
    Along with the above conductance of $P$, let us consider
$$
    \hat\BW_P(v,w) =\WW(\hat\GG_P^{vw});\quad \hat\BW_P =\sup_{v,w} \hat\BW_P(v,w).
 $$
By (\ref{T^vw}),
\begin{equation}\label{hat}
    \BW_P \leq \hat\BW_P + \frac{C}{\mu}.
\end{equation}

Take a pair of vertices, $v$ and $w$.
Let $Q^v\cap T^{vw}=\{v'\}$ and  $Q^w\cap T^{vw}= \{w'\}$.
Recall that depth of the puzzle pieces $Q^v$ and $Q^w$ is equal to $r=(2\n-1)\q+1$,
and so depends only on $\q$ and $\n$. Let $\EE^r$ be the lift of the equipotential of level $r$
to $\BP$.

For any horizontal multicurve $\gamma\in \hat\GG_P^{vw}$,
one of the following two  possibilities can occur:

\ssk\nin $\bullet$
  $\gamma$ crosses the equipotential $\EE^r$, and hence it contains an arc $\gamma'$ connecting $\EE^r$
to $T^{vw}$; By Lemma \ref{C} and the Parallel Law,
the width of this family of curves is bounded by $2^{r+\n} / \mu$ (here $2^{\n-1}$
is a bound on the number of connected components of $P\sm K_P$);

\ssk\nin $\bullet$
 contains two disjoint multicurves,
$\de^v$ and $\de^w$, that do not cross $\EE^r$ and such that $\de^v$ connects $\di^v\BP$ to $T^{vw}$,
while $\de^w$ connects $T^{vw}$ to $\di^w\BP$. Then $\de^v$ contains a multicurve that can
be lifted to a horizontal multicurve in $\BQ^v$ connecting $\di^v\BQ^v$ to $\di^{v'}\BQ^v$,
 an similarly for  $\de^w$.

\ssk
By the Series and Parallel Laws,
$$
     \hat\BW_P(v,w) \leq \BW_{Q^v}(v, v') \oplus \BW_{Q^w}(w,w')+ \frac{2^{r+\n}}{\mu}
           \leq \BW_{Q^v} \oplus \BW_{Q^w}+\frac{2^{r+\n}}{\mu}.
$$
But $\BW(Q^v)=\BW(Q^w)= \BW(P)$  since $Q^v$ and $Q^w$ are univalent pullbacks of $P$.
Hence
$$
   \BW_{Q^v}\oplus \BW_{Q^w}  \leq \frac{1}{2} \BW_P.
$$
Putting the last two estimates together and taking the supremum over all pairs of vertices $(v,w)$ of $P$,
we conclude that
$$
   \hat\BW_P\leq \frac{1}{2}\BW_P + \frac{2^{r+\n}}{\mu}.
$$
 Together with  (\ref{hat}) it yields:
$$
   \BW_P\leq \frac{1}{2} \BW_P + \frac{C}{\mu},
$$
and the conclusion follows.
\end{proof}

\begin{prop}\label{separation prop}
  $\Bd_{Y^1}(\alpha, \alpha') \geq C^{-1}\mu. $
\end{prop}

\begin{proof}
   Since the map $f^{(\n-1)\q}: \BP\ra \BY^1$ is a branched covering
that maps $\di_0\BP$ to  $\di_0\BY^1$, any curve
$\gamma\in \GG_{Y^1}^{\alpha \alpha'}$ can be lifted to a curve
$\gamma^*\in \cup\GG_P^{vw}$, where the union is taken over all pairs of vertices of $P$.
Hence
$$
 \LL (\GG_{Y^1}^{\alpha \alpha'}) \geq \bigoplus_{v,w} \LL(\GG_P^{vw})\geq \frac{1}{N\BW_P},
$$
where $N$ is the number of pairs $(v,w)$.
The conclusion follows.
\end{proof}

 Lemma \ref{observe} and Proposition~\ref{separation prop} imply:

\begin{cor}\label{mu vs Y-R}
   $\displaystyle{\frac{\mu}{2} \leq C\, \psimod (Y^0, R).}$
\end{cor}


\begin{cor}\label{slow change}
Let $f: (\BU, K)\ra (\BU, K)$ be a renormalizable $\psi$-quadratic-like map
with decoration  parameters $(\q, \n)$,
and let $f'=f^p : (\BU', K')\ra (\BU',K')$ be its first renormalization.
Then
$$
      \min\{ \mod (\BU, K), 1/2\} \leq C\,  \mod(\BU', K'),
$$
where $C=C(\q,\n)$.
\end{cor}

\begin{proof}
  This follows from Lemma \ref{simple} and Corollary \ref{mu vs Y-R}
by noticing that $(\BE^{\chi-1}, \KK) = (\BU', K')$.
\end{proof}


\subsection{Conclusion}\label{conclusion}

Everything is now prepared for the main results.
 Corollary \ref{R-L dist} and Corollary \ref{mu vs Y-R} imply:

\begin{thm}[Improving of the moduli: bounded decoration parameters]\label{mod improve}
For any parameters $\bar\q,\bar\n$ and any $\rho>0$,
 there exist  ${\underline p} \in \N$ and $\eps>0$ with the following property.
Let $f: (\BU, K)\ra (\BU, K)$ be a renormalizable $\psi$-quadratic-like map
with  decoration parameters $(\q,\n)\leq (\bar\q, \bar\n)$,
and let $f'=f^p : (\BU', K')\ra (\BU',K')$ be its first renormalization.
Then
$$
 \{ p\geq \underline p\ \ \mathrm{and}\ \mod(\BU', K')< \eps \}  \imply \mod (\BU, K) < \rho \mod(\BU', K').
$$
\end{thm}

%\begin{proof}
%It follows from  Corollary \ref{R-L dist} and Corollary \ref{mu vs Y-R}
%by noticing that $(\BE^{\chi-1}, \KK) = (\BU', K')$.
%\end{proof}





\begin{rem}  The logic of this theorem can be adjusted so that it would sound more
like an ``improvement in the future'' rather than ``worsening in the past'':

{\it  For any parameters $\bar\q, \bar\n$ of a Misuirewicz limb,
 there exists  $\underline p \in \N$ and $\eps>0$ such that
$$
    \mod(\BU', K') \geq   2\,  \mod (\BU, K)
$$
provided  $p\geq \underline p$ and $\mod (\BU, K) < \eps/2$.}
\end{rem}


\msk
  Theorem \ref{mod improve},
together with Lemma \ref{ql extension},
implies Theorem \ref{high periods thm} from the Introduction.

To derive the Main Theorem, we will combine Theorem \ref{mod improve}
with the following result  (Theorem 9.1 from \cite{K}):

\begin{thm}[Improving of the moduli: bounded period]\label{improving mod-2}
For any $\rho\in (0,1)$, there exists $\underline p=\underline p (\rho)$
such that for any $\bar p\geq \underline p$, there exists $\eps=\eps(\bar p)>0 $ with the following property.
Let $f: (\BU,K)\to (\BU,K)$ be primitively renormalizable $\psi$-quadratic-like map,
% with  $\underline p\leq p\leq \bar p$.
and let $f'=f^p = (\BU',K') \ra (\BU',K')$ be the corresponding renormalization.
Then
% Let $\gamma$ and $\gamma'$ be the geodesics associated with $f$ and $R^s f$ respectively.\\
$$
   \{\underline p\leq p\leq \bar p\ \ \mathrm{and}\ \mod(\BU'\sm K') < \eps\}   \imply \mod(\BU\sm K) < \rho \mod(\BU'\sm K').
$$
\end{thm}

\begin{rem}
  Unlike Theorem \ref{mod improve}, in Theorem \ref{improving mod-2} the map
$f'$  {\it is not necessarily the first renormalization of} $f$.
On the other hand, in Theorem  \ref{improving mod-2},
the scale $\eps$ {\it depends on the upper bound} $\bar p$,
while in Theorem \ref{mod improve} it does not.
\end{rem}


We say that an infinitely renormalizable $\psi$-ql map $f$
belongs to the {\it decoration class} $(\bar\q, \bar\n)$ if the decoration parameters $(\bar\q_n, \bar\n_n)$
of the renormalizations $R^n f$ are all bounded by $(\bar\q, \bar\n)$.

Let us now put the above two theorems together:

\begin{cor}\label{improving-3}
 For any $(\bar\q, \bar\n)$,   there exist an $\eps>0$ and $l\in \N$ with the following property.
For any  infinitely renormalizable $\psi$-ql map $f$  of decoration class $(\bar\q, \bar \n)$
with renormalizations $R^n f: (\BU_n, K_n) \ra (\BU_n, K_n)$,
if $\mod(\BU_n\sm K_n)< \eps$, $n\geq l$,  then $\mod (\BU_{n-l}\sm K_{n-l})<\mod(\BU_n\sm K_n) /2$.
\end{cor}



\begin{proof}
  Given an infinitely renormalizable $\psi$-ql map $f$ with renormalizations
$R^n f: (\BU_n, K_n)\ra (\BU_n\sm K_n)$, we let $\mu_n(f)= \mod (\BU_n, K_n).$
Arguing by contradiction, we find a sequence $f_i$ of infinitely renormalizable $\psi$-ql maps
of decoration class $(\bar\q, \bar\n)$ and sequences $\eps_i\to 0$ and $n(i) \to \infty$
such that:

\ssk\nin P1:
   $\mu_{n(i)}(f_i) <\eps_i$;

\ssk\nin P2:
  $ \mu_{n(i)}(f_i) < 2 \mu_k (f_i)$,  $k=0,1,\dots,  n(i)-1 $.
\ssk


Let $R^n f_i $ be the renormalization of $R^{n-1} f_i$ with period $p_n(f_i)$.
Applying  then the diagonal process,  we can also assume the following property:


\ssk\nin P3:
   $p_{n(i)-s}(f_i)\to \pi_s\in \N \cup \{\infty\}$ for  $s=0,1, \dots$.
\ssk

We  let $\bar s \in \Z_{\geq 0} \cup\{\infty\}$ be the first moment for which $\pi_s=\infty$
(with understanding that $\bar s=\infty$ if such a moment does not exist).

Let us consider two cases:

\ssk
{\it Case 1}: $\bar s<\infty$.
Applying consecutively Corollary \ref{slow change}, we conclude that for sufficiently big $i$,
$$
   \mu_{n(i)-s}(f_i) \leq C^s \mu_{n(i)}(f_i), \quad s\leq \bar s.
$$
Let $\rho\in (0, 1/2 C^{\bar s})$.
By Theorem \ref{mod improve}, for all sufficiently big $i$,
$$
   \mu_{n(i)-\bar s-1}(f_i) \leq \rho\,  \mu_{n(i)-\bar s}(f_i).
$$
 Putting the last two estimates together, we conclude that for all sufficiently big $i$,
$$
   \mu_{n(i)-\bar s-1}(f_i) < \frac{1}{2} \mu_{n(i)}(f_i),
$$
  contradicting assumption (P2).

\ssk
{\it Case 2}: $\bar s=\infty$.
  Take an $s$ such that
$$
   \bar p\equiv \pi_0 \pi_1\dots \pi_s > \underline p,
$$
where $\underline p = \underline p(1/2)$ comes from Theorem \ref{improving mod-2}.
By this theorem, for sufficiently big $i$,
$$
   \mu_{n(i)- s - 1}(i) < \frac{1}{2} \mu_{n(i)}(i),
$$
contradicting again assumption (P2).
\end{proof}


We are ready to prove the Main Theorem, in an important refined version.
We say that a family $\MM$ of little Mandelbrot copies (and the corresponding
renormalization combinatorics) has {\it beau}%
\footnote{According to Dennis Sullivan,
   ``beau'' stands for ``bounded and eventually universal''. }
{\it a priori} bounds
if there exists an $\eps =\eps(\MM)>0$ and a function $N: \R_+\ra \N$
with the following property.
Let $f: U\ra V$ be a quadratic-like map with $\mod(V\sm U)\geq \de>0$
that is  at least $N=N(\de)$ times renormalizable. Then for any $n\geq N$,
the $n$-fold renormalization of $f$ can be represented  by a quadratic-like map $R^n f: U_n\ra V_n$ with
$\mod(V_n\sm U_n)\geq \eps$.

\begin{bbrmt}
%\proclaim Beau Bounds (Refined Main Theorem).
  For any parameters $(\bar\q, \bar\n)$,
the family of renormalization combinatorics of decoration class $(\bar\q, \bar\n)$ has
beau a priori bounds.
\end{bbrmt}

\begin{proof}
   Let $\eps>0$ and $l$ come from Corollary \ref{improving-3},
and $C>0$ comes from  Corollary \ref{slow change}.
We will use notation $\mu_n(f)$ from the proof of Corollary \ref{improving-3}.
Assume that for some $\de>0$, there is a sequence of $\psi$-ql maps $f_i$ in question with $\mu_0(f_i)\geq \de$,
while $\mu_{n(i)}(f_i)<\eps$, where $n(i)\to \infty$. Let $n(i)=k_i l+r_i$ where $0\leq r_i<l$.
Then by  Corollaries \ref{improving-3} and \ref{slow change},
$$
   \mu_0(f_i)\leq C^l \eps /2^{r_i} \to 0\quad {\mathrm as}\ i\to \infty.
$$
This contradiction proves the beau bounds for the moduli $\mu_n(f)$ of $\psi$-ql maps.
The beau bounds for ordinary quadratic-like maps follow by Lemma \ref{ql extension}.
\end{proof}
